% !TeX root = ../papers/nc-p3.tex
\section{半素な特殊隅からの被約trace降下}
\label{sec:finite-semiprime-specialization}
\begin{theorem}\label{thm:finite-semiprime-specialization}
定理\ref{thm:finite-domain-specialization-parameters}の有限四辺の閉包上で、
「標準正則かつ一隅が整域」は「標準正則かつ一隅が半素」と同値である。
従って同所の564次までの検査は後者も判定する。検査を通らない正則点があれば
全ての隅が非零の冪零両側イデアルを持つ必要がある。
\end{theorem}
\begin{proof}[小さい極小成分の排除]
$A=e_0Be_0$ が半素で、窓は正則とする。$A$ はアフィンPI環なので
古典的商は有限個の単純Artin環の積である。ここで必要な有限性を説明する。
アフィンPI環のJacobson根基は冪零であり、半素性により零である。
全単純は有限次元で次元に一様上界がある。有限個の次元の表現多様体の単純開集合を
既約成分に分け、一般表現の核を取れば有限個の素イデアルの交わりが零となる。
一般像が係数体拡大後に全行列環を張るので各核は素である。
各素PI商はPosnerの定理でGoldieであり、こうした有限個の極小素商を持つ半素環も
Goldieである（消去イデアルの昇鎖条件と有限一様次元は有限積への埋込み、
または各極小素に対応する互いに直交した本質的なイデアルの和から従う）。
極小素は次数作用で個々に保たれる。

極小素商のうちGK次元が1以下のものは存在しない。
その成分だけに台を持つ非零な斉次元 $f$ を取る。他の極小素を全て掛けたイデアルは
指定した極小素に含まれないからである。有限次元成分なら $fB$ は有限次元で矛盾する。
実際 $fA$ はその成分の中に入り、各 $fBe_i$ は $fA$ と長さ3以下の経路で張られる。
GK次元1のアフィン素PI代数は中心上有限である（Small--Stafford--Warfield）。
従ってその次数付き中心のNoether正規化として一つの正次数元を取り、斉次な $a\in A$ に持ち上げられる。
$af=fa$ が $A$ で成立し、$fA/a fA$ は有限次元である。
左 $a$ 倍は $N=fB$ の次数 $l>0$ の右加群自己射となり単射である。
実際 $afb=0$ なら、右から頂点0に戻す全経路 $c$ に対して $afbc=0$、
成分上の正則性から $fbc=0$ となる。右から十分長い全経路が $fb$ を消すことになり、
AS消滅により $fb=0$。同じ短い末尾分解で $N/aN$ は有限次元である。
$0\to N(-l)\to N\to N/aN\to0$ に次数付きHomを $B$ へ取る。
有限次元加群からの第0,1 Extが零なので
$\underline{\operatorname{Hom}}_B(N,B)\simeq\underline{\operatorname{Hom}}_B(N(-l),B)$。
一方 $N$ は一生成で $B$ は下に有界なので、これはHom全体が零となることを意味する。
包含に反する。
\end{proof}

\begin{proof}[離散付値方向の商環と特異指標]
一般窓からこの点への曲線の局所離散付値環 $O$ を取り、$t$ をそのパラメータ、
$\mathcal A$ を零番隅の自由な次数別格子とする。
一般窓は三十節点とその局所模型の開集合に選ぶ。
$\mathcal A/t^n$ で、$A$ において正則となる元全体を左右分母として局所化する。
これらは各厚みでも正則でありOre条件を持つ。実際剰余でOre方程式
$au=sb$ を分母 $u$ で解き、一致が $t^j$ まで得られたら、
誤差を $t^j$ で割った剰余について再び $s$ に対するOre方程式を解く。
右から次の分母を掛けて $b$ を補正すれば精度が一つ上がる。反対側も同様である。
局所化の逆極限を $Q$ とすると、$Q/t$ は先の半単純古典商、$Q$ は $t$ 平坦かつ完備である。
これは各厚みの局所化の核が $t^jQ_n$ となることと
$t^jQ_n/t^{j+1}Q_n=Q_1$ から従う。
行列単位を完備環に持ち上げると $Q=\prod_v M_{r_v}(T_v)$ で、
$T_v/t$ は可除環である。異なる単純ブロック間の隅は剰余零で、
$t$ で何度でも割れるので零となる。$T_v[1/t]$ は可除環で $t$ が値1の離散付値を与える
（先頭剰余が可逆）。$\mathcal A$ は $Q$ に入り、$\mathcal A[1/t]\cap Q=\mathcal A$
という格子の飽和性を持つ。ここで交わりは $Q[1/t]$ 内で取り、
飽和性は $A\hookrightarrow Q/t$ の単射による。

一つのブロック $Q'_v=M_{r_v}(T_v[1/t])$ を固定し、一般アフィン中心のその中での像の
分数体を $F_v$ とする。中心の像は局所化と完備化後も中心に入る。
一般隅はその中心有限代数なので、$F_v$ と一般隅の像が生成する代数 $J_v$ は有限次元である。
$Q'_v$ の中に稠密であることを確かめる。正則剰余を持つ分母 $s$ はこのブロックで可逆で、
$J_v$ に作用する左乗法が単射なため有限次元性で逆も $J_v$ に入る。
これで各厚みの任意の元を近似でき、$t$ も逆にできるから稠密である。
$F_v$ の完備化を $\widehat F_v$ とすれば $Q'_v$ は $J_v\otimes\widehat F_v$ の商である。
実際、像は稠密であり完備中心体上有限次元の部分空間なので閉である。
特に $Q'_v$ はこの体上有限次元であり、$Q_v/t$ は $F_v$ の剰余体上有限次元である。
後者は剰余で一次独立な元を持ち上げ、係数の最小付値で割っても独立なことから従う。

一般窓の中心の $F_v$ 値の指標が錐の特異軌跡または頂点に入ることはない。
その集合は三十本の射影点の錐であり次元1以下なので、入れば
$F_v$ は $\operatorname{Frac}O$ 上の超越次数が高々1である。
付値の剰余体も $k$ 上の超越次数が高々1である。剰余の代数独立元の付値零の
持ち上げは、係数の最小付値を除いて還元することで $\operatorname{Frac}O$ 上独立と分かる。
よって特殊隅のこの素成分はGK次元1以下である。ここではアフィン素PI代数の
GK次元が商中心体の超越次数であることを用いた。これは一般行列のtrace環で有限次元行列表示を取り、
有限個の行列成分による上界と中心局所化の下界を比較する通常のPI次元公式である。
剰余中心が有限型と事前に分からなくても、有限生成代数の係数を収める有限生成部分体で
この上界を取れる。これは前段の排除に反する。

従って一般中心の各像指標はAzumaya軌跡に入る。$J_v$ は $F_v$ 上次数4の中心単純環であり、
完備体への拡大も単純なので $Q'_v$ と同型である。特に一般隅の被約traceはこの
ブロックの被約traceに写る。行列環では対角成分の被約traceの和であり、
離散付値を持つ可除隅の被約traceは整で剰余と両立する
（定理\ref{thm:quantum-component-intersection}と前節の付値基底の証明）。
よって一般隅の任意の格子元の被約traceは全ブロックで整であり、格子の飽和性により
再び $\mathcal A$ に属する。剰余では各単純Artin成分の次数も4で、
traceはその被約traceになる。特に特殊隅の中心元の持ち上げにこのtraceの四分の一を
施せば、剰余で元の中心元そのものになる。
\end{proof}

\begin{proof}[四隅の両立と半素性から整性へ]
一隅が半素なら $B$ も半素である。非零な両側イデアル $I$ について
$e_0 I e_0\ne0$ が成立する。右および左から頂点0に達する経路の全てで
非零元が消えれば有限次元片側イデアルが生じるからである。
よって $I$ が冪零なら $e_0 I e_0$ も非零冪零となり矛盾する。各隅も従って半素である。

一般全中心の格子の特殊化が零番隅中心の全体を与えることを示す。
斉次な $c\in Z(A)$ を持ち上げた一つの零番ループ $b$ の一般の被約traceを取る。
その四つの一般中心成分はいずれも特殊化可能である。
零番は前段による。他位置でも、経路語 $b=u v$ をその位置で巡回して $v u$ としたものの
被約traceを取れば同じ一般中心成分を得る。
一般商では四頂点が同じ行列階数を持つため、二つの同じサイズの行列の積の
traceの巡回性そのものである。正の四の倍数長の語は全位置を通り、和にも線型に使える。
他の各隅にも前段を適用しているので同時に整である。
その四分の一で $c$ は全中心に延び、零番制限は上の両側イデアルの性質から単射である。
よって $Z(B)=Z(A)$。各特殊隅はこの中心上有限である。
実際各元は剰余の被約特性係数（Newton公式）でモニックな四次式を満たし、
前節の高さ定理の議論が適用できる。各頂点間の加群は一方の隅から短い語で生成される。
これで $B$ は有限生成中心上有限となる。斉次Noether正規化の四元で窓の剰余は有限であり、
正則性と定理\ref{thm:successively-central-cutoff}によりこれら四元上自由となる。
頂点射影の整数Hilbert重複度の議論で全ての非零合成が非零となり、前節の判定に戻る。
\end{proof}
